% ==============================================================================
% ACM journal / conference template (skeleton)
% Based on acmart v2.19  —  CTAN: https://ctan.org/pkg/acmart
% Repository: https://github.com/borisveytsman/acmart
% Class acmart is contained in standard TeX Live / MiKTeX, so this file
% compiles out-of-the-box with:  pdflatex -> bibtex -> pdflatex -> pdflatex
%
% Document-class formats (pick exactly one):
%   manuscript  single-column, double-spaced, review copy
%   acmsmall    small ACM journals (default if no format given)
%   acmlarge    large ACM journals
%   acmtog      ACM Transactions on Graphics
%   sigconf     ACM conference proceedings (two-column)        <-- commonest
%   siggraph    SIGGRAPH conference
%   sigplan     SIGPLAN conferences (PLDI, POPL, ICFP, ...)
%   sigchi      SIGCHI conferences (CHI, ...)
%   acmengage   ACM EngageCSEducation
%
% Common options appended to the format:
%   review        adds line numbers + reviewer-friendly spacing
%   anonymous     double-blind anonymisation
%   author-version, nonacm, timestamp, doi, isbn, ...
%
% Typography (acmart defaults):
%   Body font  : Linux Libertine (serif) / Linux Biolinum (sans)
%   Math font  : Libertinus Math
%   Caption    : 8-9pt
%   Line spacing: single (review: double)
% Citation   : ACM Reference Format  (\bibliographystyle{ACM-Reference-Format})
%              numbered by default; author-year where the journal requires it
% Figures    : single column = \columnwidth  |  double = \textwidth (figure*)
% Captions   : below figures, above tables
% ==============================================================================
\documentclass[sigconf,review]{acmart}
% For a journal manuscript instead:
% \documentclass[manuscript,review,anonymous]{acmart}
% For SIGPLAN conferences:
% \documentclass[sigplan,review]{acmart}

% ---- Core packages: acmart already loads amsmath, hyperref, graphicx,
%       booktabs, array, xcolor, natbib, etc. Do NOT reload them. ----
% (Add only packages acmart does NOT provide, e.g. listings, minted.)

% ---- ACM metadata helpers (required for camera-ready) ----
\acmConference[Conference 'YY]{Conference Name}{Month DD--DD, 20YY}{City, Country}
% \acmISBN{978-x-xxxx-xxxx-x/XX/XX}
% \acmDOI{10.1145/XXXXXXXX.YYYYYYYY}
\acmYear{20YY}
\copyrightyear{20YY}
% \acmPrice{\$15.00}

% ---- Rights / permission (set when you receive the e-Rights form) ----
\setcopyright{none}              % use 'acmcopyright' for camera-ready
% \setcctype{cc}               % Creative Commons (some ACM journals)

% ==============================================================================
% Title, authors, affiliations
% ==============================================================================
\title{A Concise and Informative Title for the ACM Submission}

% Each author block: \author, \authornote (optional), \affiliation, \email.
\author{First Author}
\authornote{Both authors contributed equally.}
\email{first.author@univ.edu}
\affiliation{%
  \institution{University of Y}
  \department{Department of X}
  \city{City}
  \country{Country}
}

\author{Second Author}
\email{second.author@inst.edu}
\affiliation{%
  \institution{Institute of W}
  \city{City}
  \country{Country}
}

\author{Senior Author}
\email{senior.author@univ.edu}
\orcid{0000-0000-0000-0000}     % ACM encourages ORCID
\affiliation{%
  \institution{University of Y}
  \department{Department of X}
  \city{City}
  \country{Country}
}

% ---- Abstract, CCS, keywords (must come BEFORE \maketitle in acmart) ----
\begin{document}

\begin{abstract}
% TODO: 150-250 words. State the problem, the approach, the key results with
% numbers, and the broader impact. No citations, no undefined abbreviations.
\end{abstract}

% ---- ACM Computing Classification System (required for camera-ready) ----
% Generate the right codes at https://dl.acm.org/ccs, then paste the CCSXML
% block and the \ccsdesc lines. The CCSXML block is commented out here so the
% skeleton compiles without network access; uncomment it for camera-ready.
% \begin{CCSXML}
% <ccs2012>
%    <concept>
%        <concept_id>10003752.10003790.10011740</concept_id>
%        <concept_desc>Computer systems organization~Embedded systems</concept_desc>
%        <concept_significance>500</concept_significance>
%    </concept>
% </ccs2012>
% \end{CCSXML}

\ccsdesc[500]{Computer systems organization~Embedded systems}
\ccsdesc[300]{Computer systems organization~Real-time systems}

\keywords{
% TODO: 3-6 keywords separated by semicolons (ACM convention).
keyword one; keyword two; keyword three
}

\maketitle

% ==============================================================================
% 1. Introduction
% ==============================================================================
\section{Introduction}
\label{sec:intro}
% TODO: context -> gap -> approach -> contributions -> roadmap. Use \cite{}
% (ACM Reference Format numeric) or \citet{}/\citep{} if author-year is on.

% ==============================================================================
% 2. Related Work
% ==============================================================================
\section{Related Work}
\label{sec:related}
% TODO: cluster prior work and position your contribution.

% ==============================================================================
% 3. Methods
% ==============================================================================
\section{Methods}
\label{sec:methods}
% TODO: problem formulation, proposed approach, theoretical analysis.

\subsection{Problem formulation}
% TODO.

\subsection{Proposed method}
% TODO.

% ==============================================================================
% 4. Experiments
% ==============================================================================
\section{Experiments}
\label{sec:experiments}
% TODO: setup, main results, ablation, qualitative examples, statistics.

\begin{figure}[t]
\centering
% \includegraphics[width=\columnwidth]{figures/result.pdf}
\fbox{\rule{0pt}{3cm}\rule{\columnwidth}{0pt}}
\caption{Single-column figure. Caption goes BELOW the figure. Use
\texttt{\textbackslash columnwidth} for single column; for a double-column
figure use \texttt{figure*} and \texttt{\textbackslash textwidth}.}
\label{fig:result}
\end{figure}

\begin{table}[t]
\caption{Table caption goes ABOVE the table (ACM convention).}
\label{tab:results}
\centering
\begin{tabular}{lcc}
\toprule
Method   & Metric A       & Metric B       \\
\midrule
Baseline & 0.812          & 0.754          \\
Ours     & \textbf{0.891} & \textbf{0.832} \\
\bottomrule
\end{tabular}
\end{table}

% ==============================================================================
% 5. Conclusion
% ==============================================================================
\section{Conclusion}
\label{sec:conclusion}
% TODO: recap contribution and headline result; one sentence of future work.

% ==============================================================================
% Acknowledgements (use \section* to keep it out of the TOC; ACM formats it)
% ==============================================================================
% \begin{acks}
% TODO: funding, grant numbers, individuals. ACM wraps this in a styled block.
% \end{acks}

% ==============================================================================
% References  —  ACM Reference Format
% Compile:  pdflatex main  ->  bibtex main  ->  pdflatex main  ->  pdflatex main
% ==============================================================================
\bibliographystyle{ACM-Reference-Format}
\bibliography{references}
% Inline fallback (uncomment if not using BibTeX):
% \begin{thebibliography}{99}
% \bibitem{ref_example} First Author and Second Author. 20YY. Title of the
%   paper. \emph{ACM Trans. X} Y, Z (Mon. 20YY), 1--25.
% \end{thebibliography}

\end{document}
